%
% This is a borrowed LaTeX template file for lecture notes for CS267,
% Applications of Parallel Computing, UCBerkeley EECS Department.
% Now being used for CMU's 10725 Fall 2012 Optimization course
% taught by Geoff Gordon and Ryan Tibshirani.  When preparing 
% LaTeX notes for this class, please use this template.
%
% To familiarize yourself with this template, the body contains
% some examples of its use.  Look them over.  Then you can
% run LaTeX on this file.  After you have LaTeXed this file then
% you can look over the result either by printing it out with
% dvips or using xdvi. "pdflatex template.tex" should also work.
%
\documentclass[twoside]{article}
\setlength{\oddsidemargin}{0.25 in}
\setlength{\evensidemargin}{-0.25 in}
\setlength{\topmargin}{-0.6 in}
\setlength{\textwidth}{6.5 in}
\setlength{\textheight}{8.5 in}
\setlength{\headsep}{0.75 in}
\setlength{\parindent}{0 in}
\setlength{\parskip}{0.1 in}
%
% ADD PACKAGES here:
%
\usepackage{amsmath,amsfonts,graphicx}
%
% The following commands set up the lecnum (lecture number)
% counter and make various numbering schemes work relative
% to the lecture number.
%
\newcounter{lecnum}
\renewcommand{\thepage}{\thelecnum-\arabic{page}}
\renewcommand{\thesection}{\thelecnum.\arabic{section}}
\renewcommand{\theequation}{\thelecnum.\arabic{equation}}
\renewcommand{\thefigure}{\thelecnum.\arabic{figure}}
\renewcommand{\thetable}{\thelecnum.\arabic{table}}
%
% The following macro is used to generate the header.
%
\newcommand{\lecture}[4]{
   \pagestyle{myheadings}
   \thispagestyle{plain}
   \newpage
   \setcounter{lecnum}{#1}
   \setcounter{page}{1}
   \noindent
   \begin{center}
   \framebox{
      \vbox{\vspace{2mm}
    \hbox to 6.28in { {\bf Physics 467/667: Thermodynamics/Statistical Physics
	\hfill Spring 2017} }
       \vspace{4mm}
       \hbox to 6.28in { {\Large \hfill Lecture #1: #2  \hfill} }
       \vspace{2mm}
       \hbox to 6.28in { {\it Lecturer: #3 \hfill Scribes: #4} }
      \vspace{2mm}}
   }
   \end{center}
   \markboth{Lecture #1: #2}{Lecture #1: #2}
   \vspace*{4mm}
}
%
% Convention for citations is authors' initials followed by the year.
% For example, to cite a paper by Leighton and Maggs you would type
% \cite{LM89}, and to cite a paper by Strassen you would type \cite{S69}.
% (To avoid bibliography problems, for now we redefine the \cite command.)
% Also commands that create a suitable format for the reference list.
\renewcommand{\cite}[1]{[#1]}
\def\beginrefs{\begin{list}%
        {[\arabic{equation}]}{\usecounter{equation}
         \setlength{\leftmargin}{2.0truecm}\setlength{\labelsep}{0.4truecm}%
         \setlength{\labelwidth}{1.6truecm}}}
\def\endrefs{\end{list}}
\def\bibentry#1{\item[\hbox{[#1]}]}
%Use this command for a figure; it puts a figure in wherever you want it.
%usage: \fig{NUMBER}{SPACE-IN-INCHES}{CAPTION}
\newcommand{\fig}[3]{
			\vspace{#2}
			\begin{center}
			Figure \thelecnum.#1:~#3
			\end{center}
	}
% Use these for theorems, lemmas, proofs, etc.
\newtheorem{theorem}{Theorem}[lecnum]
\newtheorem{lemma}[theorem]{Lemma}
\newtheorem{proposition}[theorem]{Proposition}
\newtheorem{claim}[theorem]{Claim}
\newtheorem{corollary}[theorem]{Corollary}
\newtheorem{definition}[theorem]{Definition}
\newenvironment{proof}{{\bf Proof:}}{\hfill\rule{2mm}{2mm}}
% **** IF YOU WANT TO DEFINE ADDITIONAL MACROS FOR YOURSELF, PUT THEM HERE:
\newcommand\E{\mathbb{E}}
\begin{document}
%FILL IN THE RIGHT INFO.
%\lecture{**LECTURE-NUMBER**}{**DATE**}{**LECTURER**}{**SCRIBE**}
\lecture{4}{Energy in Thermal Physics}{Qiang Zhu}{scribe-name1,2,3}
%\footnotetext{These notes are partially based on those of Nigel Mansell.}
% **** YOUR NOTES GO HERE:
% Some general latex examples and examples making use of the
% macros follow.  
%**** IN GENERAL, BE BRIEF. LONG SCRIBE NOTES, NO MATTER HOW WELL WRITTEN,
%**** ARE NEVER READ BY ANYBODY.

\section{Some useful equations} % Don't be this informal in your notes!
\begin{equation} \label{idealgas} PV = nRT = NkT \end{equation}
\begin{equation} \label{Avogadro} N = n \times N_A \end{equation}
\begin{equation} \label{PV-micro} PV = Nm{\overline v_x^2} = NkT\end{equation}
\begin{equation} \label{eqpartition} U_{thermal} = N \cdot f \cdot \frac{1}{2}kT \end{equation}
\begin{equation} \label{1stlaw} \Delta{U} = Q + W \end{equation}
\begin{equation} \label{work3} \Delta W = - P \Delta{V} \end{equation}
Q: heat transfer by conduction, convection, radiation.\\
W: mechanical/electric/chemical work. \\

\section{Compression Work}

\begin{equation} \label{duw5} VT^{f/2} = \textrm{const}. \end{equation}
\begin{equation} \label{duw6} V^{\gamma}P = \textrm{const}. \end{equation}
where ${\gamma}$ = $\frac{f+2}{f}$ is called adiabatic exponent.\\

{\bf Exercises}\\
(Problem 1.38):\\
Two identical bubbles rist from the bottom of a lake to its surface.\\
Bubble A, rise quickly (adiabatic condition).\\
Bubble B, rise slowly (isothermal condition).\\
Which buble is larger in the end?
 \\\\\\

(Problem 1.40):\\
In Lec02, we have dertimined that 
  \begin{equation} \frac{dP}{dz} = -\frac{mg}{kT} P \end{equation}
and obtained the pressue $P$ as a function of height as follows,
  \begin{equation} P(z) = P(0) \exp(-mgz/kT) \end{equation}

1) show that $T$ and $P$ has the following relation
  \begin{equation} \frac{dT}{dP} = \frac{2}{f+2} \frac{T}{P} \end{equation}
2) find a formulat for $dT/dz$ like $dP/dz$.\\
3) estimate the temperature at the peak of Mount Everest (8848 km).\\\\\\\\\\\\\\\\\\\\\

\section{Heat Capacities}
The heat capacity of an object is the amount of heat needed to raise its temperature,
  \begin{equation} C = \frac{Q}{\Delta{T}} = \frac{\Delta{U}-W}{\Delta{T}}\end{equation}
Specific heat capacity is a more fundamental metric,
  \begin{equation} c = \frac{C}{m}. \end{equation}

There are two types of circumstances
\begin{enumerate}
\item{constant volume}, $C_V$
\item{constant pressure}, $C_P$
\end{enumerate}
Obviously, $C_P$ is more representive than $C_V$.
According to eq \ref{eqpartition}
  \begin{equation} C_V = \frac{\partial U}{\partial{T}} 
                       = \frac{\partial} {\partial{T}}(\frac {NfkT}{2})
                       = \frac{Nfk}{2}
  \end{equation}
Under constant pressure,
  \begin{equation} C_P = \frac{\partial (U-W)} {\partial{T}} 
                       = \frac{\partial} {\partial{T}}(\frac {NfkT}{2}) - \frac{\partial W}{\partial{T}}
  \end{equation}
  \begin{equation}
                \frac{\partial W}{\partial{T}} = -P\frac{\partial V}{\partial T} 
                                                = -P\frac{\partial (NkT/P)}{\partial T} 
                                                = -Nk
  \end{equation}
therefore, 
  \begin{equation} C_P = C_V + Nk. \end{equation}
using ideal gas law.\\\\

{\bf Discussions}: 
\begin{enumerate}
\item{What's the relation between $C_V$ and $C_P$ for solids?}
\item{$C_V$ of 1 mole of H$_2$ gas as a function of temperature (Figure 1.13)}
\item{$C_P$ of 1 mole of elemental solids as a function of temperature (Figure 1.14)}
\end{enumerate}

{\bf Exercises}\\
(Problem 1.44): Look up the table of thermodynamic data at room temperature. Browse through the $C_P$ values in this table, and understand them according to the equipartition theorem.\\
\begin{tabular}{|c | c | c | c |}
\hline
Type    & Examples ~~~~~~~~~~~~~~~& Ideal Value ~~& Anomaly~~~~~~~~~~~~~~~~~~~ \\\hline
monoatomic gas  &   &   &   \\\hline
diatomic gas    &   &   &   \\\hline
polyatomic(linear) gas &   &   &   \\\hline
polyatomic gas &   &   &   \\\hline
Elemental solid &   &   &   \\\hline
Binary solid &   &   &   \\\hline

\end{tabular}\\\\


\section{Enthalpy}
Constant pressure processes occur quite often in chemical recations and phase transformations. Keeping track of the work done during these processes gets to be a pain. Any idea to make it more convenient? Instead of talking about the energy, we can agree to add the work due to $PV$ in the given environment. This results a new quantity called the enthalpy ($H$),
\begin{equation} H = U + PV \end{equation}
Conveniently, we can express 
\begin{equation} C_P = \bigg (\frac {\partial H}{\partial T} \bigg)_{P}.
\end{equation}
Think about
\begin{enumerate}
\item{diamond and graphite under high pressure}
\item{formation enthalpy of H$_2$ and O$_2$ gas combine to water}
\end{enumerate}

\end{document}

